\documentclass[journal]{IEEEtran}

\usepackage[utf8]{inputenc}
\usepackage{cite}
\usepackage{amsmath,amssymb,amsfonts}
\usepackage{graphicx}
\usepackage{textcomp}
\usepackage{xcolor}
\usepackage{booktabs} 
\usepackage{multirow} 
\usepackage{tabularx}
\raggedbottom
\begin{document}

\title{Generative AI for IoT Systems: Architectures, Challenges, and Future Directions}

\author{Nada Gawpah\\
 \textit{Department of Information Technology}\\ 
 \textit{Faculty of Computer and Information Technology}\\
 \textit{Sana’a University}
\textit{ Sana’a, Yemen}\\
\textit {nada.gawbah@su.edu.ye}

}

\maketitle

\begin{abstract}
The development of the digital landscape toward 2030 offers strongly the convergence of Generative AI (GenAI) and the Internet of Things (IoT). The integration appears in the Generative IoT (GIoT). It makes a Transformation in industrial and sensory intelligence. This paper provides a comprehensive review of the integration of Large Language Models (LLMs) and generative frameworks based on IoT. This Paper discusses challenges, applications, trends and solutions in GANs based on IoT. GenAI offers a powerful capability in supporting data, RF sensing, and autonomous decision-making but deploying it is critical because of a lot of challenges such as resource constraints, catastrophic forgetting, and the "black-box" nature of deep models which have been covered in this paper. We introduce a comparative analysis between traditional and generative AI. This paper covered most applications that used in different areas like Industrial, Healthcare, Smart Cities, Multimedia, Security, and Privacy. The current findings in this review summarizes the resilient, transparent, and efficient GIoT architectures.
\end{abstract}

\begin{IEEEkeywords}
Generative AI(GansAI); Internet of Things (IoT); Large Language Models (LLMs).
\end{IEEEkeywords}

\section{Introduction}
\IEEEPARstart{O}{ver} the two decades, The Internet of Things become a very transformative concept and as we are so close to 2030. The new model appears known as Internet of Senses (IoS) is emerging. which is different from traditional virtual reality, The Internet of Senses introduce a very different multi-sensory experiences acknowledging that in our physical reality, our perception extends far beyond just sight and sound. It's consists of group of senses. This review shows that the modern techniques driving immersive multi-sensory media realize it's capabilities and application [27]. The integration of the Artificial Intelligence (AI) models with the Industrial Internet of Things (IoT) systems appears as a pivotal area of research which provides a unique chance to improve the industrial processes and operational efficiency [11]. Artificial intelligence (AI), especially deep learning proved a success in various areas like computer vision, speech recognition, and natural language processing. AI introduced into the IoT as artificial intelligence of things (AIoT) [8].
The economic and business management for the cognitive digital twin-based Internet of Robotic Things (IoRT) in flexible collaborative environments and in 3D virtual factories (from the ability of automation interoperability for sustainable business performance and innovation) developed incrementally according to 3D computer vision-based production and multisensory tracking technologies [31]. However, transferring a high amount of heterogeneous data, realizing the complex environments from these data and then making smart decisions in the appropriate time consider as difficult thing. Black-box nature for the Artificial Intelligence (AI) models do not allow the users to understand the output that created with these models. In AI applications the results are not the only important thing but also the decision paths to the results [9].
Generative Artificial Intelligence (GAI) is an important branch of Artificial Intelligence (AI) which is responsible for Content that generated by (AIGC) has recently a huge concern from the scientific community to big tech giants and governments [19]. The development of Generative AI (GenAI) and Large Language Models (LLMs) show a large advancement in artificial intelligence. It has capability to generate diverse content such texts, images, and code [28]. ChatGPT has a high global concern marking a milestone in the Generative Artificial Intelligence. While Generative AI was available during the past decade, but the introduction of ChatGPT has an effect of research and innovation in this area [6]. GAI platforms do not learn just from the entry dataset but it also depends on the interaction of data to integrate, modify, and create new information to support its learning process [19]. Generative Artificial Intelligence (GAI) has changes various areas, like machine learning, healthcare, business, and entertainment, because of its marvelous capability to create real data [18].
By the integration GAI with modern Internet of Things (IoT), Generative Internet of Things (GIoT) appears to hold a large capability to make revolution in various sides of society. This enabling the IoT applications to be more efficient and intelligent [4]. Large language models (LLMs) have showed recognizable capacities in various tasks. The integrating these capacities in the Internet of Things (IoT) applications has much research attention recently [14]. The tourism industry started a transformative revolution which has a rapid advancement and integration of digital technologies. One of these technologies generative artificial intelligence (Gen AI), natural language processing (NLP), and the Internet of Things (IoT) are proving its effects [17]. The Generative AI (GenAI) has capacities in overcome limitations in traditional Radio Frequency (RF) sensing methods in the Internet of Things (IoT) systems. GenAI reviews state-of-the-art GenAI techniques and their application to RF sensing problems.The GenAI's has an ability to generate high quality artificial data, improving signal quality, and integrate multi-media data. That will provide a strong solution to reconstruct RF environment, locations, and imaging. GenAI also improves the efficiency and performance of IoT devices by enabling it to adapt to new environments and visible tasks [22].
Because of the rapid expansion in the Internet of Things (IoT) devices, the health care sector becomes responsible in a huge amount of real-time data in the power systems (RES), which provides a motivation for custom health metrics [10]. The fast integration Required for renewable energy sources (RES) in the modern energy systems a huge challenge in quality accuracy, grid stability, and energy improvement. Generative Artificial Intelligence (Gen-AI) and architectures like Generative Adversarial Networks (GANs), Variational Autoencoders (VAEs), and transformers introduce a new capability to overcome data scarcity, non-linearity, and uncertainty in RES-dominated systems [13]. Accurate forecasting of water considered an important thing to refining the management and operation of water distribution networks (WDNs)which ensures sustainable use of resource and service reliability [12].
The Internet of Things (IoT) has a huge growth, which provide a good connectivity and convenience. However, this growth caused significant cybersecurity concerns because of the heterogeneity of IoT devices. And the wide deployment, and inherent computational limitations. A lot of IoT devices has the prioritize for the low cost and user simplicity for the security. That will make them available for breaches and exploitation. including data breaches, Distributed Denial-of-Service (DDoS) attacks, and malware infections. The resource mostly prevents constraints of resources in some IoT devices from adopting contemporary encryption and authentication techniques, the long age and widespread deployment of these devices make the managing security upgrades and patches difficult.
Machine Learning (ML) introduces as an advancing technology that shows a huge promise to solve IoT security issues. It effects widely and improves research in cyber threat detection by distinguish patterns, identifying anomalies, and predicting possible threats in real-time. This capability of the pre-emptive response to vulnerabilities and threats. ML algorithms help to determine abnormal activities, flag potential security risks for the early detection of cyberattacks. In addition, improving signature-based detection by automating signature creation and updates. Furthermore, ML can analyze a huge amount of data from IoT devices to distinguish hidden connections, identify trends, and expecting future threats, which enabling to take an informed security decisions and resource allocation efficiently [23].
Artificial Intelligence (AI) introduces a powerful base of techniques to prosses these challenges. [10]. Although there are a lot of chances introduced by the Generative AI to improve the IoT systems Performance, but this integration faces a complex challenge such as resource constraints, latency, privacy issues, and data security. This paper aims to provide a holistic view of IoT-based GenAI architectures, and a comparative analysis with traditional and Generative AI. In addition, this paper presents a research issues and optimization techniques, critical findings and outlines future research directions.\\
In this paper, a review of Generative Artificial Intelligence for IoT Systems is conducted, and it is organized as follows:
\textbf{Section II:} Architectures and Paradigms of IoT-Based GenAI.\\
\textbf{Section III:} Comparative Analysis: Traditional AI vs. Generative AI in IoT.\\
\textbf{Section IV:} Key Research Issues and Challenges\\
\textbf{Section V:} Key Solutions and Optimization Techniques\\
\textbf{Section VI:} Applications and Use Cases\\
\textbf{Section VII:} Critical Findings and Lessons Learned\\
\textbf{Section VIII:} Future Trends and Research Directions
\begin{figure}[htbp]
    \centering
    \includegraphics[width=1.1\linewidth]{nn.jpg} 
    \caption{Structural mind map of Generative AI Integration in IoT Ecosystems }
    \label{fig:my_label}
\end{figure}
\begin{figure}[htbp]
    \centering
    \includegraphics[width=1.1\linewidth]{mm.jpg} 
    \caption{Conceptual architecture for the review paper}
    \label{fig:my_label}
\end{figure}
\subsection{Research Methodology}
This review introduced a comparative and methodology to evaluate the role of Gen-AI predicting the renewable energy and system optimization. The literature selection involved a systematic search in the high effected journals indexed in Scopus covering the period from 2020 to 2025. The key criteria for selecting the models and applications are:\\
\textbf{•Relevance:} Models and applications that can be applicable to IIoT scenarios and in areas such as predictive maintenance, quality control, supply chain management, and energy improvement.\\
\textbf{•Impact:} Research and case studies shows an improvement in efficiency, accuracy, or productivity because of applying deep learning in the IIoT.\\
\textbf{•Novelty:} Inclusion of well-established models such as (CNNs, RNNs, LSTMs) and emerging models (GANs, autoencoders) to provide a wide perspective on trends and innovations.\\
\textbf{•Publication Quality:} choosing reviewed articles, high-impact conference papers to ensure the reliability and validity of the information presented.
\begin{table}[htbp]
\caption{Inclusion and Exclusion Criteria}
\label{table:criteria}
\centering
% تغيير lll إلى ccc للتوسيط
\begin{tabular}{@{}ccc@{}} 
\toprule
\textbf{Criterion} & \textbf{Inclusion Criteria} & \textbf{Exclusion Criteria} \\ \midrule
Time Period & \begin{tabular}[c]{@{}c@{}}Recent studies published \\ between 2020 and 2026\end{tabular} & \begin{tabular}[c]{@{}c@{}}Studies published \\ before 2020\end{tabular} \\ \addlinespace
Study Type & Journals & \begin{tabular}[c]{@{}c@{}}Conferences, articles, \\ and short papers\end{tabular} \\ \addlinespace
Core Theme & \begin{tabular}[c]{@{}c@{}}Generative AI (GenAI, LLMs, \\ GANs) with IoT\end{tabular} & \begin{tabular}[c]{@{}c@{}}Traditional AI/ML \\ without GANs\end{tabular} \\ \addlinespace
Focus Area & \begin{tabular}[c]{@{}c@{}}Security, Privacy, \\ Optimization, and Resources\end{tabular} & \begin{tabular}[c]{@{}c@{}}General AI studies \\ without IoT context\end{tabular} \\ \addlinespace
Language & Full-text English & Other languages \\ \bottomrule
\end{tabular}
\end{table}
\begin{table}[htbp]
\caption{PRISMA SELECTION PROCESS SUMMARY}
\label{table:prisma_results}
\centering
\begin{tabular}{@{}ccc@{}}
\toprule
\textbf{Stage} & \textbf{Description} & \textbf{Count} \\ \midrule
Identification & \begin{tabular}[c]{@{}c@{}}Initial records identified via IEEE, \\ ScienceDirect, ACM, and Springer\end{tabular} & 150 \\ \addlinespace
Screening & \begin{tabular}[c]{@{}c@{}}Records after removing duplicates \\ and non-English papers\end{tabular} & 124 \\ \addlinespace
Initial Filtering & \begin{tabular}[c]{@{}c@{}}Records screened based on \\ Title and Abstract relevance\end{tabular} & 45 \\ \addlinespace
Final Inclusion & \begin{tabular}[c]{@{}c@{}}Final studies included in the \\ systematic review\end{tabular} & 31 \\ \bottomrule
\end{tabular}
\end{table}
\begin{table}[htbp]
\caption{DISTRIBUTION OF SELECTED REFERENCES}
\label{table:distribution}
\centering
\begin{tabular}{@{}ccc@{}}
\toprule
\textbf{Research Domain} & \textbf{Reference Numbers} & \textbf{Percentage} \\ \midrule
\begin{tabular}[c]{@{}c@{}}Fundamentals \&\\ Architectures\end{tabular} & [4, 5, 8, 11, 27] & 16\% \\ \addlinespace
\begin{tabular}[c]{@{}c@{}}Security, Privacy \&\\ Trust\end{tabular} & \begin{tabular}[c]{@{}c@{}}[1, 2, 3, 9, 16, \\ 18, 23, 24, 28]\end{tabular} & 29\% \\ \addlinespace
\begin{tabular}[c]{@{}c@{}}Resource \& Hardware\\ Optimization\end{tabular} & \begin{tabular}[c]{@{}c@{}}[14, 15, 22, 25, \\ 30]\end{tabular} & 16\% \\ \addlinespace
\begin{tabular}[c]{@{}c@{}}Domain-Specific\\ Applications\end{tabular} & \begin{tabular}[c]{@{}c@{}}[7, 10, 12, 13, 17, \\ 20, 21, 26, 29, 31]\end{tabular} & 32\% \\ \addlinespace
\begin{tabular}[c]{@{}c@{}}Generative Models\\ (GANs, Transformers, LLMs)\end{tabular} & [6, 19] & 7\% \\ \bottomrule
\end{tabular}
\end{table}
\subsection {Perspectives of AI in IoT}
This paper presents the (AIoT) architecture Then, it presents a progress in AI research for IoT. Advancements in Generative AI provided AIoT with chances to have advantages from generative models such as large language models (LLMs) to perceive, interpret, and present IoT sensor data [15]. Studies describe that the integration of generative language models such as ChatGPT and Gemini with IoT devices contributes in interaction with processes and physical elements in many ways [26].
\section{ARCHITECTURES AND PARADIGMS OF IOT-BASED GENAI}
The Architectures That used is multi-agent architecture with two special large LLMs each of them has a specific task to support the accuracy of responses. especially that LLMs may not perform good in dealing with long text chain or tasks that required long planning. The architecture of proposed model integrates IoT with health monitoring, advanced pre-processing, feature extraction, and generative AI to access to a special healthcare visions [10].
\subsection{Taxonomy of Generative Models}
Gen AI consist of a group of technics that designed to generate a new content of text and images to music and codes according learned patterns and data. The backbone of Gen AI is in neural networks, especially deep neural networks (DNNs) [17]. Deep generative models (DGMs) are classified to several categories:\\
\textbf{•Variational Autoencoders (VAEs)}\\
\textbf{•Generative Adversarial Networks (GANs)}\\
\textbf{•Autoregressive Models}\\
\textbf{•Flow-based Models}\\
\textbf{•Diffusion Models. [18, 25].}\\
A basic architecture of GAI are three fundamental components VAE, diffusion model, and GAN [19].\\
\textbf{•	Autoencoders and Variational Autoencoders (VAEs)}\\
Autoencoder is a model of neural network for the unsupervised machine learning that encrypt the input data using encoder into a low dimensional representation (encoding). Then used a decoder to decode it back to its real form (decoding). In addition, it reduces the fault of reconstruction. This model has designed to reduce Dimension, Feature Extraction, deleting noise from images, Anomaly Detection and compromising Missing Value [31, 11]. An autoencoder is consist of four components:\\
\textbf{Encoder:} Compresses input data to lower dimensions (code).\\
\textbf{Code/Bottleneck:} A layer contains the compressed representation.\\
\textbf{Decoder:} Reconstructs the code layer.\\
\textbf{Reconstruction Loss:} Measures un similarities.\\
Autoencoders consider not generative, but the Variational Autoencoder (VAE) is the replacing generative [6]. VAE models learn an encoder and decoder making them perfect for large dataset such as radio maps and personalizing health metrics [15, 10, 22].\\
\textbf{• Generative Adversarial Networks (GANs)}
GANs consist of as Goodfellow et al., introduce it two neural networks (the generator and the discriminator). They are a part of a zero competitive game [17, 11]. The generator creates realistic data from zero. On the other hand, the discriminator distinguishes and evaluates them against the real data [11, 19]. In IoT, GANs used to create real synthetic RF data, supporting signal quality, and enable robust anomaly detection [22]. And there is a lot of Various versions such as Robust GAN (RGAN) and Vehicle Synthesis GAN (VS-GAN) [20].\\
\textbf{• Diffusion Models (DMs)}
DMs is using Markov chain to adding random noise to data (forward diffusion) then omitting that noise it to generate required samples (reverse diffusion) [22]. Some papers proposed a way to generate a network traffic sample that leverages from the Conditional Denoising Diffusion Probabilistic Model (CDDPM) to fixes the issues of limited sample diversity [2]. GDMs also work as a basic for secure incentive mechanism frameworks in GIoT [4].\\
\textbf{• Transformer Models and Large Language Models (LLMs)}
"Attention Is All You Need" is a paper that introduced the transformer model that used self-attention to process chain of data [6, 17]. These models were the base of models like GPT and BERT.\\
Generative Pretrained Transformer (GPT): LLM that depend on A transformer.it generates texts that are similar to human text [6].
Architectures: The most popular is the architecture based on decoder-only such as0 (GPT-3, LaMDA). Include the Recent open-source models like LLaMa, Mistral, and Mixtral. Mixtral follows the Mixture of Experts architecture (MoE) [27].
In IoT systems, LLMs facilitate advanced multi-modal data analysis by integrating RF data with text and audio inputs for holistic vision [22].\\
\textbf{• Hybrid and Emergent Frameworks}
A hybrid generative framework combines between the GANs networks for realistic pattern and VAE models to learn organized representations [10]. These models synthesize data using GANs, transformers, and diffusion models to develop industrial worlds that is based on digital twin and production management systems [31]. Against the traditional ways like SMOTE, these generative models create a completely new distributed data that captured complex attack patterns in a better way and edge cases in cybersecurity [16].

\subsection{Synthetic Data Generation for IoT Augmentation}
In IoT scenarios, collecting data faces some logistical, financial, and privacy constraints. In addition to real world environments complexity. Also collecting Large wide attacks data and classifying it costly and needs time. Therefore, NIDS for IoT based on (Few-Shot Learning) are crucial [2]. Generative models can create an artificial attacks data to support training datasets. That will improve the diversity and robustness of models. Also, discriminative models’ success in classification and prediction. By combining them, we can achieve to high accurate predictions. including generative adversarial networks (GANs) to support data [3].
Generative AI (GAI) can leverage from the pre-training models to create new content by sitting the model parameters according to the user inputs like prompts. Also, GAI can use learning algorithms to automate content creation from existing data [4]. This support does not only facilitate the depth analysis of communication and sensing data. but also generates a huge amount of data to train the communication and sensing model. that may be difficult for the existing AI models [5]. GAI can leverage from the advanced AI algorithms to create virous, and valuable data and manipulating and modifying it [21]. Deep generative models (DGMs)used to approximate and generate a joint distribution for the training, and target data to generate similar samples to the real data, which has physical process plausibility in IIoT applications [25].
Environment mapping and fault diagnosis algorithms can help in creating artificial data across the smart robotic environments. It’s a result of the usefulness of multi-sensory extended reality and simulation modeling techniques [31]. By leveraging AI models to generate an artificial high accuracy data. it’s possible to support the limited data in the real-world which enhance the training and performance of RF signal processing algorithms. Generative AI can improve the model robustness by creating an artificial data from the learned distributions. Also involving controlled contrast, and enabling unseen circumstances. GenAI can produce a specific data for some scenarios under an environmental conditions, interference patterns, and contrasted device configurations that simulate virous settings like urban, indoor, and rural areas [22].
In GAI frameworks VAEs generated realistic and high-quality data. And that is making them in a high valuable in limited classified data cases [19]. The generator starts to generate an artificial data such as images or sounds, from random set of input data (often referred to as noise vectors) It aims to be not capable to distinguishable from realistic data. On the other hand, the discriminator is working as a judge trying to distinguish between real and the fake data. Training for GAN network is happening in a repeated way until it reaches to a balancing point where the generator learned to create vary realistic data [11]. GAIN, which works based on generative adversarial network (GAN) produces realistic imputations via adversarial learning. Compared to KNN and MICE. GAIN extracts the complex non-linear dependency of the variables and interactions [12].
Furthermore, GPT could be used to generate an artificial text data, which can then be transferred through GAN network to generate realistic images. This integration can be functioned to generate more artificial data for computing vision models, audio, security, and text models, which can improve its accuracy [24]. Generative AI introduce a lot of advantages dealing with incomplete and spared datasets. This is too useful for renewable energy systems when the historical data may be insufficient or unavailable [13].
Generation artificial data and Privacy-preserving is a too important for the risk-aware testing. the models based on transformer reaches accuracy levels up to 92\% for the ICMP protocols. while Conditional VAE (CVAE)-based methods combine differential privacy with generative modeling to produce un known datasets for intrusion detection [1]. PPT such as differential privacy federated learning and secure multi-party computation is used to generate an artificial data or making computations and preserving privacy at the same time during the data training [18].
Most of the articles (89\%) use generative AI to improve the performance of the IDS. while, a few models (11\%) think from the cyber attackers view about how to generate attacks that can deceive detection. To improve intrusion detection, we can classify the uses of GenAI into 2 classes: data augmentation including class balancing, and data reconstruction [16]. Artificial data can also simulate hard environmental circumstances in industrial IoT to predict faulty equipment and optimizing maintenance schedules. Similarly, in healthcare IoT, Artificial RF data can enhance model training for patient monitoring systems [22, 20]. Finally, DL models generate high-level abstraction and implementable visons that provide feedback to support IoT systems and their services [30].
\subsection{Smart City and IoT Architectures}
The environment of the smart city has an advanced communication techniques and an IoT infrastructure, which requires specific communication technologies and an advance systems architecture to adopt all the applications of smart cities [7]. Bringing AI closer to the edge provide a great solution for achieving high performance and improving the quality of service (QoS) and quality of experience (QoE) for sensitive applications of delaying [30]. Compared to the traditional IoT, the computing component of AIoT focused on compute tasks that is AI-oriented. also, the sensing, networking, and communication components become enable by AI. that allows billions of daily devices to attack the modern AI [15].\\
\section{COMPARATIVE ANALYSIS: TRADITIONAL AI VS. GENERATIVE AI IN IOT}
Traditional AI/ML models, such as LSTM and CNN still valuable due to their light computational requirements and its interpretability. specially in predicting tasks. These models perform well under stable circumstances, but it effected in high variability environments and data-scarce settings [13]. Traditional data analysis methods, such as CNN-based and U-Net based methods, depend on a large amount of categorized data and face difficulties in dealing with high-dimensional data [21]. Before 2014, all the deep learning models was descriptive, focusing on summarizing current data patterns. However, the introducing of GANs in 2014 made the generative models learn the underlying probability distribution of the data instead of just creating mapping abstractions [6, 19].
\begin{table}[h]
\centering
\caption{Comparative of Traditional AI / ML and Generative AI (GenAI)}
\label{tab:ai_comparison}
\setlength{\tabcolsep}{3pt} % تقليل المسافات بين الأعمدة للسماح بمساحة أكبر للنص
\small % تصغير حجم الخط لضمان عدم خروج الجدول
\begin{tabular}{@{} p{1.6cm} p{2.8cm} p{2.8cm} c @{}}
\toprule
\textbf{Feature} & \textbf{Traditional AI / ML} & \textbf{Generative AI (GenAI)} & \textbf{Ref} \\
\midrule
Data Req. & Requires less data to train & Requires large amounts of data & [24,25] \\
\addlinespace
Scalability & Less computationally demanding & Computationally demanding & [24,25] \\
\addlinespace
Core Function & Classification and regression & Modeling complex distributions & [24,25] \\
\addlinespace
Interaction & Rigid and specific & Natural language and reasoning & [24,25] \\
\bottomrule
\end{tabular}
\end{table}
\subsection{Key Differences and Advantages}
The following points clarify the shift from traditional to generative approaches in IoT:

\begin{itemize}
    \item \textbf{Data Handling and Augmentation:} Traditional methods depend on large, clean datasets, but GenAI can generate high-quality artificial data to support the limited data [32]. Unlike traditional sampling techniques like SMOTE, generative models produce a new data distribution that can capture complex attack patterns [16].
    
    \item \textbf{Robustness to Noise:} Traditional methods have very limited robustness against varying conditions, while GenAI can learn the underlying distribution of both signal and noise, which allows it to generalize and reconstruct incomplete information [32, 22].
    
    \item \textbf{Generalization and Adaptation:} Traditional models struggle with static network structures and adapting to dynamic environments, but GenAI can generalize to IoT devices and adapt to new environments and tasks [32, 22].
    
    \item \textbf{Security and Intrusion Detection:} Traditional intrusion detection systems (IDS) most of the time rely on rules or specified signatures, while modern dynamic challenges like APTs and zero-day vulnerabilities require complex and adaptive mechanisms provided by GenAI [16, 1].
\end{itemize}

\section{Key Research Issues and Challenges}
The development of Generative AI introduces several chances and challenges [6]. It provides users with perfect solutions, but there are some limits that should be considered for efficient implementation in IoT environments.

\subsection{Data and Model Complexity}
\begin{itemize}
    \item \textbf{Data Complexity:} IoT devices generate data that is huge in volume, variety, and velocity, making it necessary to provide them with efficient and scalable algorithms.
    
    \item \textbf{Distributional Bias:} Relying on artificial data generated by GANs and VAEs may be effective for rare condition augmentation, but it may introduce distributional bias if the synthetic samples do not fully capture real-world variability [10].
    
    \item \textbf{Reliability of Unsupervised Learning:} Unsupervised learning methods can face problems in identifying the importance of results and require manual analysis [23].
    
    \item \textbf{Catastrophic Forgetting:} Transfer learning models can suffer from catastrophic forgetting when updated frequently with small datasets of patients [10].
\end{itemize}

\subsection{Hardware and Architectural Constraints}
\begin{itemize}
    \item \textbf{Resource-Intensive Nature of DL:} Deep learning methods are considered resource-intensive, representing an issue for IoT devices [23].
    
    \item \textbf{Hardware Constraints:} Hardware constraints on edge devices can limit the implementation of complex models and the processing of high-frequency physiological data in real-time [10].
    
    \item \textbf{Device Heterogeneity and Interoperability:} The diversity in IoT devices—varied architectures, protocols, and operating systems—makes it difficult to develop a globally applicable solution. Integration with heterogeneous IoT devices poses interoperability challenges in maintaining consistent data quality across varied sensor specifications [10, 23].
\end{itemize}
\subsection{Computational and Resource Constraints}
Internet of Robotic Things (IoRT) digital twins operate through cloud computing and big data infrastructures, yet they involve complex computational tasks, often relying on deep reinforcement learning. However, many IoT devices prioritize low cost and user simplicity over security, which leads to significant resource constraints. Deploying generative models such as Transformers and diffusion models is highly challenging because lightweight hardware often lacks the necessary processing power and memory, raising concerns regarding real-time, low-latency responses. Even lightweight Deep Neural Networks (DNNs) specialized for object detection contain millions of parameters, making it infeasible to run multi-stage AIoT control algorithms on platforms with limited memory. Furthermore, hosting Large Language Models (LLMs) on mobile units like UAVs is often avoided due to their massive size, high computing demands, and energy consumption, which reduce operational flight time. For enterprises that avoid commercial LLM services due to security reasons, deploying open-source LLMs locally remains a challenge because they possess limited capacities.

\subsection{Accounting Constraints}
AI systems require large amounts of data, meaning that processing this data can consume significant resources. Retraining or fine-tuning LLMs is considered costly and resource-intensive, especially because these models can quickly become out-of-date. Therefore, optimizing computing and memory resources for embedded devices is essential for efficient implementation.

\subsection{Security, Privacy, Ethical Implications, and Trust}
The vulnerabilities inherent in AIoT pose security concerns, as resource-constrained devices struggle to implement high-level security measures. Furthermore, privacy is a primary concern, as the collection of location, healthcare records, and behavior patterns by AIoT can raise risk levels. Additionally, reaching final results without interpretation may result in intentional or unintentional discrimination or trust-related issues.

\begin{itemize}
    \item \textbf{Dynamic Nature of Threats:} Cyber threats evolve continually, requiring continuous learning and adaptation in threat detection methods. This dynamic nature often necessitates manual intervention, which acts as a limiting factor.
    \item \textbf{Vulnerability of ML Models:} Adversaries may exploit vulnerabilities in ML models, making it necessary to develop robust methods to counter adversarial attacks and manipulation.
    \item \textbf{Explainability and Bias:} Reaching results without any explanation may result in intentional or unintentional discrimination or trust issues.
    \item \textbf{Implementation Challenges:} There are significant challenges in managing computational complexity, ensuring data security, and addressing ethical issues during real-world implementation.
    \item \textbf{Data Privacy:} GenAI models often require access to sensitive data, such as energy usage or Electronic Health Records (EHR). Without strong protection, these systems risk exposing sensitive infrastructure vulnerabilities.
    \item \textbf{Ethical Concerns:} Using GenAI raises concerns related to fake videos, misinformation, and the creation of deepfakes, which can violate individual rights regarding their images and voices.
    \item \textbf{Algorithmic Bias:} Bias may occur in federated learning due to heterogeneous data distributions. Many DL models trained on general data are biased against vulnerable demographics such as children and the elderly.
\end{itemize}

\subsection{Adversarial Attacks}
Although GenAI models are used for vulnerability detection, they remain vulnerable to adversarial manipulations. Techniques like Adaptive Gradient based on Word Embedding Perturbations (AG-WEP) can deceive text classifiers. Adversaries may also add imperceptible noise to deceive models, causing them to generate deceptive images or "deepfakes," which impact many domains ranging from forensic analysis to autonomous vehicles.
\section{Key Solutions and Optimization Techniques}
Addressing the aforementioned challenges requires several optimization and architectural strategies, as detailed below:

\subsection{Model Compression and Efficiency}
\begin{itemize}
    \item \textbf{Quantization \& Pruning:} Utilizing fixed-point representations, such as INT8 for weights and activations, reduces computational complexity and bandwidth. Furthermore, pruning methodologies streamline model complexity \cite{11, 16}.
    \item \textbf{Knowledge Distillation:} This technique reduces model size and inference time while maintaining acceptable detection accuracy \cite{16}.
    \item \textbf{Lightweight Models:} Employing lightweight models tailored for specific tasks minimizes device load and makes edge deployment more feasible \cite{22}.
\end{itemize}

\subsection{Architectural Solutions}
\begin{itemize}
    \item \textbf{Edge and Fog Computing:} Moving data processing closer to the source (Edge) reduces latency and network congestion \cite{30, 31}. A Cloud-Edge-End cooperative storage scheme can strategically partition sensory data to improve performance and security \cite{15}.
    \item \textbf{Federated Learning (FL):} FL enables model training with privacy preservation across distributed devices without the need to share raw data \cite{13}.
    \item \textbf{Blockchain Integration:} Utilizing blockchain secures GIoT management through the sharing of decentralized production data \cite{31, 4}.
\end{itemize}

\subsection{Explainability and Robustness (XAI)}
\begin{itemize}
    \item \textbf{XAI Methods:} Applying methods such as SHAP and LIME addresses the "black-box" problem by providing clear explanations for AI decisions and improving transparency \cite{2, 9, 3}.
    \item \textbf{Resilient Training:} Implementing adversarial training—where models are exposed to adversarial examples during the training phase—helps protect models against data poisoning and prompt injection \cite{1}.
\end{itemize}

\subsection{Data Enhancement and Imputation}
\begin{itemize}
    \item \textbf{Data Imputation:} Transformers, such as BERT, can process missing sensor readings by understanding contextual relationships \cite{22}.
    \item \textbf{Super-resolution:} Diffusion models are highly efficient in generating data from sparse latent spaces to reconstruct incomplete information \cite{22}.
    \item \textbf{Transfer Learning:} This enables rapid model adaptation to new parameters without complete retraining, ensuring interoperability with various hardware \cite{10}.
\end{itemize}
\begin{table}[ht]
\centering
\caption{Challenges of GANs and Solutions}
\label{table:gans_challenges}
\small
\setlength{\tabcolsep}{3pt} % تقليل المسافة الجانبية بين الأعمدة لتوفير مساحة
\begin{tabular}{@{} p{1.6cm} p{2.4cm} p{2.8cm} c @{}}
\toprule
\textbf{Category} & \textbf{Key Challenge} & \textbf{Proposed Solution} & \textbf{Ref} \\
\midrule
Data Complexity & High volume/velocity & Scalable algorithms \& Compression & [10, 11, 16] \\
\addlinespace
Model Stability & Catastrophic forgetting & Transfer Learning & [10] \\
\addlinespace
Resource Cons. & High memory/power demands & Distillation \& Lightweight models & [16, 22, 23] \\
\addlinespace
Hardware Limits & Edge processing inability & Quantization (INT8) \& Pruning & [11, 16] \\
\addlinespace
Latency & Cloud dependency & Edge and Fog Computing & [30, 31, 22] \\
\addlinespace
Interoperability & Device heterogeneity & Transfer Learning \& Cloud-Edge-End & [10, 15] \\
\addlinespace
Privacy/Security & Sensitive data risks & Federated Learning \& Blockchain & [13, 31, 4] \\
\addlinespace
Robustness & Adversarial vulnerability & Resilient Training & [1] \\
\addlinespace
Trust \& Ethics & Black-box decisions & Explainable AI (XAI) & [2, 3, 9] \\
\addlinespace
Data Integrity & Missing readings & Data Imputation \& Super-resolution & [22] \\
\bottomrule
\end{tabular}
\end{table}
\section{Applications and Use Cases}
The integration of Generative AI (GenAI) with IoT has pioneered a new paradigm for interacting with the environment, revolutionizing various verticals such as industry, healthcare, education, and smart city infrastructure \cite{26, 27}.

\subsection{Smart Manufacturing and Industrial IoT (IIoT)}
In the Industrial Internet of Things (IIoT), AI analyzes data from connected devices to facilitate intelligent operational decision-making.

\begin{itemize}
    \item \textbf{Swarm Robotics and IoRT:} Mechanisms of swarm robotic coordination and mobile IoRT devices are monitored to enable seamless cooperation within the immersive industrial metaverse \cite{31}.
    \item \textbf{Predictive Maintenance:} Unlike traditional methods limited to static data, GenAI analyzes real-time sensor data—such as temperature, vibration, and pressure—to identify warning signals of imminent failures, thereby extending system lifespan \cite{11, 23}.
    \item \textbf{Process Optimization:} AI identifies inefficiencies in manufacturing lines by analyzing real-time data, suggesting improvements to reduce waste and optimize processing times \cite{11}.
    \item \textbf{Structural Health Monitoring:} Diffusion Models (DMs) can convert sparse RF input data into detailed structural maps to monitor critical infrastructure like bridges and pipelines \cite{22}.
\end{itemize}

\subsection{Healthcare and Biomedical Technologies}
GenAI-based systems are facilitating a transition toward personalized medicine and advanced remote care:

\begin{itemize}
    \item \textbf{Personalized Monitoring:} Frameworks utilizing wearable IoT devices enable improved diagnoses for chronic conditions, such as diabetes, through continuous monitoring and personalized care plans \cite{10, 22}.
    \item \textbf{Telemedicine and Assistive Tech:} AIoT systems facilitate remote diagnosis, on-site disease detection, and assistive technologies for individuals with disabilities \cite{15, 20}.
    \item \textbf{Medical Imaging:} GANs are employed to generate rare pathological images for training diagnostic algorithms, while VAEs are used to simulate potential failure components in medical devices \cite{13}.
    \item \textbf{Elderly Care:} Environmental sensors in smart homes provide context-aware monitoring, issuing safety alerts and early warnings for individuals suffering from cognitive impairments \cite{20}.
\end{itemize}
\section{Applications and Use Cases}
\subsection{Smart Cities and Infrastructure}
Advanced communication techniques and IoT infrastructure enable several smart city applications:
\begin{itemize}
    \item \textbf{Traffic Flow and Transportation:} GenAI predicts vehicle movements and analyzes traffic flows by reconstructing missing RF data, effectively reducing congestion \cite{22}.
    \item \textbf{Environmental Monitoring:} Denoising diffusion models (e.g., SegDiff) facilitate barrier detection and environment reconstruction, which is critical for urban planning \cite{22}.
    \item \textbf{Energy and Smart Grids:} Integrating generative architectures with energy systems enhances simulation realism, enabling grid stability analysis and energy dispatch optimization \cite{13, 30}.
\end{itemize}

\subsection{Multimedia and Cognitive Applications}
\begin{itemize}
    \item \textbf{Vision-based Applications:} GDMs and GANs allow IoT devices to interpret visual information for 3D tracking in AR, VR, and MR environments \cite{4, 15}.
    \item \textbf{Text and Voice Interaction:} LLMs like ChatGPT-4 process real-time IoT textual data, allowing users to control devices via natural language \cite{4, 26}.
    \item \textbf{Content Generation:} Specialized tools enable automated generation of images from text, code for IoT programming, and music compositions \cite{6}.
\end{itemize}

\subsection{Security and Privacy Applications}
\begin{itemize}
    \item \textbf{Vulnerability Detection:} GenAI utilizes LLMs and GAN payloads for automated fuzzing and vulnerability detection in APIs and smart contracts \cite{1}.
    \item \textbf{Anomaly Detection:} ML models monitor network traffic in real-time to identify abnormal patterns and tailor security protocols to unique IoT network features \cite{23, 25}.
    \item \textbf{Anonymization:} ADS-GAN generates artificial data approximating original distributions (e.g., EHR), preserving data utility while ensuring privacy \cite{21}.
\end{itemize}

\section{Critical Findings and Lessons Learned}
The integration of GenAI and IoT yields high-performance metrics across domains. 
\subsection{Performance Benchmarks and Efficiency}
\begin{itemize}
    \item \textbf{Resource Optimization:} GenAI-based schemes show nearly 52\% higher edge server utility compared to traditional Deep Reinforcement Learning (DRL) \cite{4}.
    \item \textbf{Sensing Accuracy:} Experimental results in ISAC physical layers show that Artificial Signal Generation (SSG) achieves a Mean Square Error (MSE) of $\approx 1.03^{\circ}$ in Direction of Arrival (DoA) estimation \cite{5}.
    \item \textbf{Cybersecurity Robustness:} Transformer-based security models for cyber threat-hunting demonstrate detection accuracy reaching 95\% \cite{24}.
\end{itemize}

\subsection{Data and Architectural Lessons}
\begin{itemize}
    \item \textbf{Data Imputation:} Models achieve a coefficient of determination $R^2$ up to 0.90, maintaining high performance even with 90\% data loss \cite{12}.
    \item \textbf{Poisoning Resilience:} SplitFed architectures are more robust to poisoning attacks than traditional Federated Learning (FL) due to smaller model portions \cite{15}.
    \item \textbf{Holistic Strategies:} No single solution addresses all IoT security gaps; a holistic strategy combining diverse techniques is essential \cite{23}.
\end{itemize}

\section{Future Trends and Research Directions}
Generative AI represents the frontier of the Fifth Industrial Revolution (5IR) \cite{6, 15}:

\begin{itemize}
    \item \textbf{Hybrid Physics-Data Models:} Future efforts must fuse physics-based simulations with GenAI pattern recognition to optimize mission-critical applications \cite{3, 13}.
    \item \textbf{Foundation Models for Sensing:} Task-agnostic models, such as Meta-Transformers, offer potential for spectrum management and beam selection without task-specific retraining \cite{22}.
    \item \textbf{Governance and Compliance:} Research must integrate GDPR-compliant risk classification directly into GAI-IoT frameworks \cite{1, 23}.
    \item \textbf{6G and Next-Generation Sensing:} ISAC systems enhanced by GAI are essential for 6G, facilitating 3D coverage and ultra-low latency \cite{5, 30}.
\end{itemize}




\section{Conclusion}
The integration of Generative AI and the IoT framework considered as start to a new space for   more intuitive and intelligent of connectivity. Although GansAI has a strong ability to completely Disrupted from healthcare to autonomous flight. This development and deployment have a black side represented in heterogeneity and resource limitations of edge devices.
This review highlights that the path to the future is in the balance between performance and transparency. Solutions Presented in this paper such as Knowledge Distillation and Quantization are essential to make the complex model appliable on lightweight hardware. On the other hand, Explainable AI (XAI) and Adversarial Training are a necessity for building trust in sensitive AIoT applications. As we are moving from simple data collection to complex content generation at the edge. The focus of future research must focus to local deployment for energy-efficient and collaborative learning based on privacy. This paper focuses in GansAI and IoT because it considered as the future of the digital technology.


\begin{thebibliography}{00}

\bibitem{1} Mirtaheri SL, et al. "Cybersecurity in the age of generative AI: A systematic taxonomy of AI-powered vulnerability assessment and risk management," \textit{Future Generation Computer Systems}, vol. 175, 2026.
\bibitem{2} Zhang Z, et al. "Trustworthy Generative Few-Shot Learning-Based Intrusion Detection Method in Internet of Things," \textit{IEEE Trans. Consum. Electron.}, vol. 71, no. 1, pp. 1992–2002, 2025.
\bibitem{3} Ankalaki S, et al. "Cyber Attack Prediction: From Traditional Machine Learning to Generative Artificial Intelligence," \textit{IEEE Access}, vol. 13, pp. 44662–44706, 2025.
\bibitem{4} Wen J, et al. "From Generative AI to Generative Internet of Things: Fundamentals, Framework, and Outlooks," \textit{IEEE Access}, vol. 12, pp. 63820-63836, 2024.
\bibitem{5} Wang J, et al. "Generative AI for Integrated Sensing and Communication: Insights from the Physical Layer Perspective," \textit{arXiv preprint arXiv:2310.01036}, 2023.
\bibitem{6} Bengesi S, et al. "Advancements in Generative AI: A Comprehensive Review of GANs, GPT, Autoencoders, Diffusion Model, and Transformers," \textit{IEEE Access}, 2024.
\bibitem{7} Alahi MEE, et al. "Integration of IoT-Enabled Technologies and Artificial Intelligence (AI) for Smart City Scenario: Recent Advancements and Future Trends," \textit{Sensors}, vol. 23, no. 11, 2023.
\bibitem{8} Zhang J, Tao D. "Empowering Things With Intelligence: A Survey of the Progress, Challenges, and Opportunities in Artificial Intelligence of Things," \textit{IEEE Internet Things J.}, vol. 8, no. 10, pp. 7789-7817, 2021.
\bibitem{9} Kök İ, et al. "Explainable Artificial Intelligence (XAI) for Internet of Things: A Survey," \textit{IEEE Internet Things J.}, vol. 10, no. 16, pp. 14764-14779, 2023.
\bibitem{10} Mahajan RA, et al. "Enhancing personalization in IoT-based health monitoring via generative AI and transfer learning," \textit{Egyptian Informatics Journal}, vol. 32, 2025.
\bibitem{11} Dini P, et al. "Overview of AI-Models and Tools in Embedded IIoT Applications," \textit{Electronics}, vol. 13, no. 12, 2024.
\bibitem{12} Velayudhan NK, et al. "Generative AI for spatio-temporal multivariate imputation and demand prediction in water distribution systems," \textit{Results in Engineering}, vol. 27, 2025.
\bibitem{13} Erdiwansyah, et al. "Emerging role of generative AI in renewable energy forecasting and system optimization," \textit{Sustainable Chemistry for Climate Action}, 2025.
\bibitem{14} Xiao B, et al. "Efficient Prompting for LLM-based Generative Internet of Things," \textit{arXiv:2502.08886}, 2025.
\bibitem{15} Siam SI, et al. "Artificial Intelligence of Things: A Survey," \textit{ACM Transactions on Sensor Networks}, 2025.
\bibitem{16} Deng Z, et al. "Generative AI in Intrusion Detection Systems for Internet of Things: A Systematic Literature Review," \textit{IEEE Open J. Commun. Soc.}, vol. 6, pp. 4689–4717, 2025.
\bibitem{17} Suanpang P, Pothipassa P. "Integrating Generative AI and IoT for Sustainable Smart Tourism Destinations," \textit{Sustainability}, vol. 16, no. 17, 2024.
\bibitem{18} Golda A, et al. "Privacy and Security Concerns in Generative AI: A Comprehensive Survey," \textit{IEEE Access}, vol. 12, pp. 48126–48144, 2024.
\bibitem{19} Vu TH, et al. "Applications of Generative AI (GAI) for Mobile and Wireless Networking: A Survey," \textit{IEEE Internet Things J.}, vol. 12, no. 2, 2025.
\bibitem{20} Naseer F, et al. "Integrating generative adversarial networks with IoT for adaptive AI-powered personalized elderly care in smart homes," \textit{Front. Artif. Intell.}, 2025.
\bibitem{21} Chen J, et al. "Generative AI-Driven Human Digital Twin in IoT-Healthcare: A Comprehensive Survey," \textit{IEEE Internet Things J.}, vol. 11, no. 1, 2024.
\bibitem{22} Wang L, et al. "Generative AI for RF Sensing in IoT Systems," \textit{IEEE Internet Things Mag.}, vol. 8, no. 2, pp. 112-120, 2025.
\bibitem{23} Alwahedi F, et al. "Machine learning techniques for IoT security: Current research and future vision with generative AI and large language models," \textit{Internet of Things and Cyber-Physical Systems}, 2024.
\bibitem{24} Ferrag MA, et al. "Generative AI for Cyber Threat-Hunting in 6G-enabled IoT Networks," \textit{Proc. IEEE/ACM CCGridW}, 2023.
\bibitem{25} De S, et al. "Deep Generative Models in the Industrial Internet of Things: A Survey," \textit{IEEE Trans. Ind. Inform.}, vol. 18, no. 9, pp. 5728-5743, 2022.
\bibitem{26} Yauri R, Espino R. "Generative Language Model Technology Integrated into an IoT Device for the Development of a Voice Assistant," \textit{WSEAS Trans. Syst.}, vol. 23, 2024.
\bibitem{27} Sehad N, et al. "Generative AI for Immersive Communication: The Next Frontier in Internet-of-Senses Through 6G," \textit{IEEE Commun. Mag.}, vol. 63, no. 2, 2025.
\bibitem{28} Aung YL, et al. "Generative AI for Internet of Things Security: Challenges and Opportunities," \textit{arXiv:2502.08886}, 2025.
\bibitem{29} Mikołajewska E, et al. "Generative AI in AI-Based Digital Twins for Fault Diagnosis for Predictive Maintenance," \textit{Applied Sciences}, 2025.
\bibitem{30} Bourechak A, et al. "At the Confluence of Artificial Intelligence and Edge Computing in IoT-Based Applications," \textit{Sensors}, 2023.
\bibitem{31} Lăzăroiu G, et al. "Cognitive digital twin-based Internet of Robotic Things... in the immersive industrial metaverse," \textit{Equilibrium}, vol. 19, no. 3, 2024.

\end{thebibliography}

\end{document}
